\section{}

\paragraph{907年}\eventanchor{LIAO-0907-EVENT-026}{LIAO-0907-EVENT-026}　耶律阿保机取得汗位，开始重组部族权力。此项以《辽史》〈本纪〉、〈营卫志〉、〈百官志〉；《宋史》辽国相关传；宋人使辽记录 为基本来源线索；编年与官修材料能够确认年序和制度用语，但对地方执行、人数、动机或社会后果的判断仍须同文书、碑刻、考古及对方叙事互校。

\paragraph{907年}\eventanchor{LIAO-0907-REGNAL-YEAR}{LIAO-0907-REGNAL-YEAR}　契丹／辽以辽太祖在位第1年作为诏令、赋役与外交文书的纪年框架。此项以《辽史》〈本纪〉、〈营卫志〉、〈百官志〉；《宋史》辽国相关传；宋人使辽记录 为基本来源线索；编年与官修材料能够确认年序和制度用语，但对地方执行、人数、动机或社会后果的判断仍须同文书、碑刻、考古及对方叙事互校。

\paragraph{908年}\eventanchor{LIAO-0908-REGNAL-YEAR}{LIAO-0908-REGNAL-YEAR}　契丹／辽以辽太祖在位第2年作为诏令、赋役与外交文书的纪年框架。此项以《辽史》〈本纪〉、〈营卫志〉、〈百官志〉；《宋史》辽国相关传；宋人使辽记录 为基本来源线索；编年与官修材料能够确认年序和制度用语，但对地方执行、人数、动机或社会后果的判断仍须同文书、碑刻、考古及对方叙事互校。

\paragraph{909年}\eventanchor{LIAO-0909-REGNAL-YEAR}{LIAO-0909-REGNAL-YEAR}　契丹／辽以辽太祖在位第3年作为诏令、赋役与外交文书的纪年框架。此项以《辽史》〈本纪〉、〈营卫志〉、〈百官志〉；《宋史》辽国相关传；宋人使辽记录 为基本来源线索；编年与官修材料能够确认年序和制度用语，但对地方执行、人数、动机或社会后果的判断仍须同文书、碑刻、考古及对方叙事互校。

\paragraph{910年}\eventanchor{LIAO-0910-REGNAL-YEAR}{LIAO-0910-REGNAL-YEAR}　契丹／辽以辽太祖在位第4年作为诏令、赋役与外交文书的纪年框架。此项以《辽史》〈本纪〉、〈营卫志〉、〈百官志〉；《宋史》辽国相关传；宋人使辽记录 为基本来源线索；编年与官修材料能够确认年序和制度用语，但对地方执行、人数、动机或社会后果的判断仍须同文书、碑刻、考古及对方叙事互校。

\paragraph{911年}\eventanchor{LIAO-0911-REGNAL-YEAR}{LIAO-0911-REGNAL-YEAR}　契丹／辽以辽太祖在位第5年作为诏令、赋役与外交文书的纪年框架。此项以《辽史》〈本纪〉、〈营卫志〉、〈百官志〉；《宋史》辽国相关传；宋人使辽记录 为基本来源线索；编年与官修材料能够确认年序和制度用语，但对地方执行、人数、动机或社会后果的判断仍须同文书、碑刻、考古及对方叙事互校。

\paragraph{912年}\eventanchor{LIAO-0912-REGNAL-YEAR}{LIAO-0912-REGNAL-YEAR}　契丹／辽以辽太祖在位第6年作为诏令、赋役与外交文书的纪年框架。此项以《辽史》〈本纪〉、〈营卫志〉、〈百官志〉；《宋史》辽国相关传；宋人使辽记录 为基本来源线索；编年与官修材料能够确认年序和制度用语，但对地方执行、人数、动机或社会后果的判断仍须同文书、碑刻、考古及对方叙事互校。

\paragraph{913年}\eventanchor{LIAO-0913-REGNAL-YEAR}{LIAO-0913-REGNAL-YEAR}　契丹／辽以辽太祖在位第7年作为诏令、赋役与外交文书的纪年框架。此项以《辽史》〈本纪〉、〈营卫志〉、〈百官志〉；《宋史》辽国相关传；宋人使辽记录 为基本来源线索；编年与官修材料能够确认年序和制度用语，但对地方执行、人数、动机或社会后果的判断仍须同文书、碑刻、考古及对方叙事互校。

\paragraph{914年}\eventanchor{LIAO-0914-REGNAL-YEAR}{LIAO-0914-REGNAL-YEAR}　契丹／辽以辽太祖在位第8年作为诏令、赋役与外交文书的纪年框架。此项以《辽史》〈本纪〉、〈营卫志〉、〈百官志〉；《宋史》辽国相关传；宋人使辽记录 为基本来源线索；编年与官修材料能够确认年序和制度用语，但对地方执行、人数、动机或社会后果的判断仍须同文书、碑刻、考古及对方叙事互校。

\paragraph{915年}\eventanchor{LIAO-0915-REGNAL-YEAR}{LIAO-0915-REGNAL-YEAR}　契丹／辽以辽太祖在位第9年作为诏令、赋役与外交文书的纪年框架。此项以《辽史》〈本纪〉、〈营卫志〉、〈百官志〉；《宋史》辽国相关传；宋人使辽记录 为基本来源线索；编年与官修材料能够确认年序和制度用语，但对地方执行、人数、动机或社会后果的判断仍须同文书、碑刻、考古及对方叙事互校。

\paragraph{916年}\eventanchor{LIAO-0916-EVENT-027}{LIAO-0916-EVENT-027}　阿保机称帝并建国号契丹。此项以《辽史》〈本纪〉、〈营卫志〉、〈百官志〉；《宋史》辽国相关传；宋人使辽记录 为基本来源线索；编年与官修材料能够确认年序和制度用语，但对地方执行、人数、动机或社会后果的判断仍须同文书、碑刻、考古及对方叙事互校。

\paragraph{916年}\eventanchor{LIAO-0916-REGNAL-YEAR}{LIAO-0916-REGNAL-YEAR}　契丹／辽以辽太祖在位第10年作为诏令、赋役与外交文书的纪年框架。此项以《辽史》〈本纪〉、〈营卫志〉、〈百官志〉；《宋史》辽国相关传；宋人使辽记录 为基本来源线索；编年与官修材料能够确认年序和制度用语，但对地方执行、人数、动机或社会后果的判断仍须同文书、碑刻、考古及对方叙事互校。

\paragraph{917年}\eventanchor{LIAO-0917-REGNAL-YEAR}{LIAO-0917-REGNAL-YEAR}　契丹／辽以辽太祖在位第11年作为诏令、赋役与外交文书的纪年框架。此项以《辽史》〈本纪〉、〈营卫志〉、〈百官志〉；《宋史》辽国相关传；宋人使辽记录 为基本来源线索；编年与官修材料能够确认年序和制度用语，但对地方执行、人数、动机或社会后果的判断仍须同文书、碑刻、考古及对方叙事互校。

\paragraph{918年}\eventanchor{LIAO-0918-REGNAL-YEAR}{LIAO-0918-REGNAL-YEAR}　契丹／辽以辽太祖在位第12年作为诏令、赋役与外交文书的纪年框架。此项以《辽史》〈本纪〉、〈营卫志〉、〈百官志〉；《宋史》辽国相关传；宋人使辽记录 为基本来源线索；编年与官修材料能够确认年序和制度用语，但对地方执行、人数、动机或社会后果的判断仍须同文书、碑刻、考古及对方叙事互校。

\paragraph{919年}\eventanchor{LIAO-0919-REGNAL-YEAR}{LIAO-0919-REGNAL-YEAR}　契丹／辽以辽太祖在位第13年作为诏令、赋役与外交文书的纪年框架。此项以《辽史》〈本纪〉、〈营卫志〉、〈百官志〉；《宋史》辽国相关传；宋人使辽记录 为基本来源线索；编年与官修材料能够确认年序和制度用语，但对地方执行、人数、动机或社会后果的判断仍须同文书、碑刻、考古及对方叙事互校。

\paragraph{920年}\eventanchor{LIAO-0920-REGNAL-YEAR}{LIAO-0920-REGNAL-YEAR}　契丹／辽以辽太祖在位第14年作为诏令、赋役与外交文书的纪年框架。此项以《辽史》〈本纪〉、〈营卫志〉、〈百官志〉；《宋史》辽国相关传；宋人使辽记录 为基本来源线索；编年与官修材料能够确认年序和制度用语，但对地方执行、人数、动机或社会后果的判断仍须同文书、碑刻、考古及对方叙事互校。

\paragraph{921年}\eventanchor{LIAO-0921-REGNAL-YEAR}{LIAO-0921-REGNAL-YEAR}　契丹／辽以辽太祖在位第15年作为诏令、赋役与外交文书的纪年框架。此项以《辽史》〈本纪〉、〈营卫志〉、〈百官志〉；《宋史》辽国相关传；宋人使辽记录 为基本来源线索；编年与官修材料能够确认年序和制度用语，但对地方执行、人数、动机或社会后果的判断仍须同文书、碑刻、考古及对方叙事互校。

\paragraph{922年}\eventanchor{LIAO-0922-REGNAL-YEAR}{LIAO-0922-REGNAL-YEAR}　契丹／辽以辽太祖在位第16年作为诏令、赋役与外交文书的纪年框架。此项以《辽史》〈本纪〉、〈营卫志〉、〈百官志〉；《宋史》辽国相关传；宋人使辽记录 为基本来源线索；编年与官修材料能够确认年序和制度用语，但对地方执行、人数、动机或社会后果的判断仍须同文书、碑刻、考古及对方叙事互校。

\paragraph{923年}\eventanchor{LIAO-0923-REGNAL-YEAR}{LIAO-0923-REGNAL-YEAR}　契丹／辽以辽太祖在位第17年作为诏令、赋役与外交文书的纪年框架。此项以《辽史》〈本纪〉、〈营卫志〉、〈百官志〉；《宋史》辽国相关传；宋人使辽记录 为基本来源线索；编年与官修材料能够确认年序和制度用语，但对地方执行、人数、动机或社会后果的判断仍须同文书、碑刻、考古及对方叙事互校。

\paragraph{924年}\eventanchor{LIAO-0924-REGNAL-YEAR}{LIAO-0924-REGNAL-YEAR}　契丹／辽以辽太祖在位第18年作为诏令、赋役与外交文书的纪年框架。此项以《辽史》〈本纪〉、〈营卫志〉、〈百官志〉；《宋史》辽国相关传；宋人使辽记录 为基本来源线索；编年与官修材料能够确认年序和制度用语，但对地方执行、人数、动机或社会后果的判断仍须同文书、碑刻、考古及对方叙事互校。

\paragraph{925年}\eventanchor{LIAO-0925-REGNAL-YEAR}{LIAO-0925-REGNAL-YEAR}　契丹／辽以辽太祖在位第19年作为诏令、赋役与外交文书的纪年框架。此项以《辽史》〈本纪〉、〈营卫志〉、〈百官志〉；《宋史》辽国相关传；宋人使辽记录 为基本来源线索；编年与官修材料能够确认年序和制度用语，但对地方执行、人数、动机或社会后果的判断仍须同文书、碑刻、考古及对方叙事互校。

\paragraph{926年}\eventanchor{LIAO-0926-EVENT-028}{LIAO-0926-EVENT-028}　辽军攻灭渤海，设置东丹等安排。此项以《辽史》〈本纪〉、〈营卫志〉、〈百官志〉；《宋史》辽国相关传；宋人使辽记录 为基本来源线索；编年与官修材料能够确认年序和制度用语，但对地方执行、人数、动机或社会后果的判断仍须同文书、碑刻、考古及对方叙事互校。

\paragraph{926年}\eventanchor{LIAO-0926-REGNAL-YEAR}{LIAO-0926-REGNAL-YEAR}　契丹／辽以辽太祖在位第20年作为诏令、赋役与外交文书的纪年框架。此项以《辽史》〈本纪〉、〈营卫志〉、〈百官志〉；《宋史》辽国相关传；宋人使辽记录 为基本来源线索；编年与官修材料能够确认年序和制度用语，但对地方执行、人数、动机或社会后果的判断仍须同文书、碑刻、考古及对方叙事互校。

\paragraph{927年}\eventanchor{LIAO-0927-EVENT-029}{LIAO-0927-EVENT-029}　耶律德光即位为辽太宗。此项以《辽史》〈本纪〉、〈营卫志〉、〈百官志〉；《宋史》辽国相关传；宋人使辽记录 为基本来源线索；编年与官修材料能够确认年序和制度用语，但对地方执行、人数、动机或社会后果的判断仍须同文书、碑刻、考古及对方叙事互校。

\paragraph{927年}\eventanchor{LIAO-0927-REGNAL-YEAR}{LIAO-0927-REGNAL-YEAR}　辽以辽太宗在位第1年作为诏令、赋役与外交文书的纪年框架。此项以《辽史》〈本纪〉、〈营卫志〉、〈百官志〉；《宋史》辽国相关传；宋人使辽记录 为基本来源线索；编年与官修材料能够确认年序和制度用语，但对地方执行、人数、动机或社会后果的判断仍须同文书、碑刻、考古及对方叙事互校。

\paragraph{928年}\eventanchor{LIAO-0928-REGNAL-YEAR}{LIAO-0928-REGNAL-YEAR}　辽以辽太宗在位第2年作为诏令、赋役与外交文书的纪年框架。此项以《辽史》〈本纪〉、〈营卫志〉、〈百官志〉；《宋史》辽国相关传；宋人使辽记录 为基本来源线索；编年与官修材料能够确认年序和制度用语，但对地方执行、人数、动机或社会后果的判断仍须同文书、碑刻、考古及对方叙事互校。

\paragraph{929年}\eventanchor{LIAO-0929-REGNAL-YEAR}{LIAO-0929-REGNAL-YEAR}　辽以辽太宗在位第3年作为诏令、赋役与外交文书的纪年框架。此项以《辽史》〈本纪〉、〈营卫志〉、〈百官志〉；《宋史》辽国相关传；宋人使辽记录 为基本来源线索；编年与官修材料能够确认年序和制度用语，但对地方执行、人数、动机或社会后果的判断仍须同文书、碑刻、考古及对方叙事互校。

\paragraph{930年}\eventanchor{LIAO-0930-REGNAL-YEAR}{LIAO-0930-REGNAL-YEAR}　辽以辽太宗在位第4年作为诏令、赋役与外交文书的纪年框架。此项以《辽史》〈本纪〉、〈营卫志〉、〈百官志〉；《宋史》辽国相关传；宋人使辽记录 为基本来源线索；编年与官修材料能够确认年序和制度用语，但对地方执行、人数、动机或社会后果的判断仍须同文书、碑刻、考古及对方叙事互校。

\paragraph{931年}\eventanchor{LIAO-0931-REGNAL-YEAR}{LIAO-0931-REGNAL-YEAR}　辽以辽太宗在位第5年作为诏令、赋役与外交文书的纪年框架。此项以《辽史》〈本纪〉、〈营卫志〉、〈百官志〉；《宋史》辽国相关传；宋人使辽记录 为基本来源线索；编年与官修材料能够确认年序和制度用语，但对地方执行、人数、动机或社会后果的判断仍须同文书、碑刻、考古及对方叙事互校。

\paragraph{932年}\eventanchor{LIAO-0932-REGNAL-YEAR}{LIAO-0932-REGNAL-YEAR}　辽以辽太宗在位第6年作为诏令、赋役与外交文书的纪年框架。此项以《辽史》〈本纪〉、〈营卫志〉、〈百官志〉；《宋史》辽国相关传；宋人使辽记录 为基本来源线索；编年与官修材料能够确认年序和制度用语，但对地方执行、人数、动机或社会后果的判断仍须同文书、碑刻、考古及对方叙事互校。

\paragraph{933年}\eventanchor{LIAO-0933-REGNAL-YEAR}{LIAO-0933-REGNAL-YEAR}　辽以辽太宗在位第7年作为诏令、赋役与外交文书的纪年框架。此项以《辽史》〈本纪〉、〈营卫志〉、〈百官志〉；《宋史》辽国相关传；宋人使辽记录 为基本来源线索；编年与官修材料能够确认年序和制度用语，但对地方执行、人数、动机或社会后果的判断仍须同文书、碑刻、考古及对方叙事互校。

\paragraph{934年}\eventanchor{LIAO-0934-REGNAL-YEAR}{LIAO-0934-REGNAL-YEAR}　辽以辽太宗在位第8年作为诏令、赋役与外交文书的纪年框架。此项以《辽史》〈本纪〉、〈营卫志〉、〈百官志〉；《宋史》辽国相关传；宋人使辽记录 为基本来源线索；编年与官修材料能够确认年序和制度用语，但对地方执行、人数、动机或社会后果的判断仍须同文书、碑刻、考古及对方叙事互校。

\paragraph{935年}\eventanchor{LIAO-0935-REGNAL-YEAR}{LIAO-0935-REGNAL-YEAR}　辽以辽太宗在位第9年作为诏令、赋役与外交文书的纪年框架。此项以《辽史》〈本纪〉、〈营卫志〉、〈百官志〉；《宋史》辽国相关传；宋人使辽记录 为基本来源线索；编年与官修材料能够确认年序和制度用语，但对地方执行、人数、动机或社会后果的判断仍须同文书、碑刻、考古及对方叙事互校。

\paragraph{936年}\eventanchor{LIAO-0936-EVENT-030}{LIAO-0936-EVENT-030}　辽支持石敬瑭建立后晋，十六州问题进入北方政治结构。此项以《辽史》〈本纪〉、〈营卫志〉、〈百官志〉；《宋史》辽国相关传；宋人使辽记录 为基本来源线索；编年与官修材料能够确认年序和制度用语，但对地方执行、人数、动机或社会后果的判断仍须同文书、碑刻、考古及对方叙事互校。

\paragraph{936年}\eventanchor{LIAO-0936-REGNAL-YEAR}{LIAO-0936-REGNAL-YEAR}　辽以辽太宗在位第10年作为诏令、赋役与外交文书的纪年框架。此项以《辽史》〈本纪〉、〈营卫志〉、〈百官志〉；《宋史》辽国相关传；宋人使辽记录 为基本来源线索；编年与官修材料能够确认年序和制度用语，但对地方执行、人数、动机或社会后果的判断仍须同文书、碑刻、考古及对方叙事互校。

